\section{Introduction}
\label{sec:intro}

\begin{figure*}[t]
    \centering
    \includegraphics[width=\linewidth]{figures/fig1_banner.pdf}
    \caption{We introduce Language-Action Pre-training (LAP), a general VLA pre-training recipe that represents low-level actions directly in natural language to supervise a vision–language backbone, and instantiate it as \textsc{LAP-3B}, the first VLA to demonstrate strong \emph{zero-shot transfer} to novel embodiments. Compared to state-of-the-art VLAs, LAP-3B learns more generalizable embodiment representations and exhibits favorable scaling behavior.
        }
    \label{fig:teaser}
    % \vspace{-1pt}
    
\end{figure*}

The ultimate promise of robot foundation policies is the ability to control \emph{any} robot embodiment to perform human-level intelligent tasks. Vision–Language–Action models (VLAs) have made remarkable progress toward this goal, leveraging pre-trained Vision–Language models (VLMs) to obtain strong visual–semantic priors for broad scene and even task generalization \cite{intelligence2025pi05visionlanguageactionmodelopenworld, lee2025molmoactactionreasoningmodels, geminiroboticsteam2025geminiroboticsbringingai, nvidia2025gr00tn1openfoundation, zheng2025xvlasoftpromptedtransformerscalable, goyal2025vla0buildingstateoftheartvlas, kim2024openvlaopensourcevisionlanguageactionmodel, kim2025finetuningvisionlanguageactionmodelsoptimizing, liu2025rdt1bdiffusionfoundationmodel}. However, despite large-scale multi-embodiment training \cite{embodimentcollaboration2025openxembodimentroboticlearning, khazatsky2025droidlargescaleinthewildrobot, agibotworldcontributors2025agibotworldcolosseolargescale, galaxea2025, wu2025robomindbenchmarkmultiembodimentintelligence}, state-of-the-art VLAs still rarely function \emph{zero-shot on new robots} with even minor differences to training embodiments (e.g., an altered gripper or wrist camera location), as shown in Fig.~\ref{fig:teaser}. In practice, deploying a VLA on a new robot often requires extensive per-embodiment adaptation, which is costly and undermines the goal of a single broadly deployable foundation model.

At its core, zero-shot cross-embodiment transfer is difficult because the embodiment space is vast and fundamentally expensive to cover with data. While task and environment diversity can scale via in-the-wild robot data \cite{chi2024universalmanipulationinterfaceinthewild, khazatsky2025droidlargescaleinthewildrobot, tao2025dexwilddexteroushumaninteractions} and internet-scale egocentric video \cite{grauman2022ego4dworld3000hours, liu2024tacobenchmarkinggeneralizablebimanual, banerjee2025hot3dhandobjecttracking, wang2023holoassistegocentrichumaninteraction, hoque2025egodexlearningdexterousmanipulation, damen2020epickitchensdatasetcollectionchallenges}, embodiment diversity remains constrained by hardware availability and data-collection logistics; for example, the Open X-Embodiment dataset aggregates data from only a few dozen robot types \cite{embodimentcollaboration2025openxembodimentroboticlearning}. Achieving zero-shot transfer therefore requires extracting maximal generalization from limited embodiment diversity, rather than relying on exhaustive coverage.

Our key observation is that zero-shot cross-embodiment transfer depends critically on how we adapt a pre-trained VLM for motor control. Concretely, the model must acquire precise control capabilities while retaining the semantic and visual representations that enable transfer across embodiments. When this balance is mismanaged, embodiment-specific action learning can override the very features needed for generalization. However, because VLMs are not pre-trained to interpret or generate motor-level, high-frequency control signals, standard VLA training can induce distributional mismatch and degrade pre-trained knowledge \cite{hancock2025actionslanguagefinetuningvlms, zhou-etal-2025-chatvla, driess2025knowledgeinsulatingvisionlanguageactionmodels}.

Motivated by this observation, we propose \emph{Language-Action Pre-training (LAP)}, a simple pre-training recipe for VLAs that adapts a pre-trained VLM backbone for motor control by training it to predict \emph{actions described in natural language}, herein referred to as \emph{language-actions}. Language-actions express the net effect of a raw end-effector action (e.g., ``move left 5 cm'', ``tilt forward 45 degrees'') in the same modality the VLM is trained to model, as shown in Fig.\ref{fig:teaser}. They can be parsed from raw actions under a \emph{fixed} coordinate convention and template without additional data annotation. 

This representation choice directly targets the distributional mismatch that arises in standard VLA training recipes, which fine-tune VLMs to predict \emph{raw actions}—continuous joint or end-effector commands, or discretizations thereof—represented by arbitrary or abstract tokens~\cite{kim2024openvlaopensourcevisionlanguageactionmodel, pertsch2025fastefficientactiontokenization, lee2025molmoactactionreasoningmodels}. Such representations carry little semantic structure that could unify behaviors across embodiments \cite{zheng2025universalactionsenhancedembodied, ye2025latentactionpretrainingvideos, bu2025univlalearningacttaskcentric, chen2025embodimentequivariantvisionlanguageactionpolicy}. Notably, this limitation persists even when a unified end-effector action space is adopted: existing VLAs still exhibit little to no zero-shot generalization to novel embodiments~\cite{kim2024openvlaopensourcevisionlanguageactionmodel, zheng2025xvlasoftpromptedtransformerscalable} (see Section~\ref{subsec:zero_shot}). Together, these results suggest that cross-embodiment transfer cannot be achieved by action space unification alone, but instead requires a fundamentally different action representation. Language-actions provide such a representation by keeping the supervision signal close to the VLM’s pre-training distribution \cite{hancock2025actionslanguagefinetuningvlms}, promoting semantically grounded features that transfer across embodiments.

With this representation in place, we next consider how to instantiate LAP in a practical VLA architecture. In principle, LAP can be applied to any VLM-based VLA by treating action generation as a language modeling problem. However, such formulations typically rely on auto-regressive language-action token sampling at inference time, which is inefficient and ill-suited for robot control \cite{kim2024openvlaopensourcevisionlanguageactionmodel, pertsch2025fastefficientactiontokenization, driess2025knowledgeinsulatingvisionlanguageactionmodels, hancock2025actionslanguagefinetuningvlms}. We therefore adopt a hybrid design that combines a LAP-trained VLM backbone with a lightweight diffusion-based action expert for continuous control, following $\pi_{0.5}$~\cite{intelligence2025pi05visionlanguageactionmodelopenworld}. This Mixture-of-Transformers architecture \cite{liang2024mixture} preserves the representation benefits of language-action supervision while enabling efficient execution.

\noindent\textbf{Statement of contributions.} We introduce Language-Action Pre-training (LAP), a general pre-training recipe for Vision--Language--Action models that represents actions as natural language to enable cross-embodiment generalization. We instantiate LAP in \textsc{LAP-3B}, which pairs a VLM backbone with a lightweight action expert for efficient real-time control. Trained on existing open-sourced robot datasets \cite{embodimentcollaboration2025openxembodimentroboticlearning, khazatsky2025droidlargescaleinthewildrobot, lee2025molmoactactionreasoningmodels}, \textsc{LAP-3B} achieves substantial zero-shot generalization to previously unseen robot embodiments, delivering roughly a 2× relative improvement (approximately 30\% absolute) over prior pre-training recipes based on alternative action representations. \emph{To the best of our knowledge, this is the first VLA system to operate zero-shot on novel real robot embodiments.}\footnote{All code, checkpoints, and data configurations are open-sourced at \href{https://lap-vla.github.io}{https://lap-vla.github.io}.}
In addition, \textsc{LAP-3B} can be fine-tuned to new embodiments and more complex dexterous tasks using significantly fewer demonstrations and gradient steps, matching prior performance with up to $2.5\times$ less data. Finally, we analyze why LAP learns more transferable embodiment representations and exhibits more stable training dynamics, demonstrate its favorable scaling behavior with increasing model size, and show that the language-action interface naturally supports co-training with auxiliary VQA objectives (e.g., motion prediction), yielding further performance gains.
